\section{射影簇之间的同构}\label{sec:isomorphisms-of-projective-varieties}

已经建立了适用于射影簇之间映射的概念, 现在转来讨论同构. 首先给出自然的定义.

\begin{Definition}{射影簇之间的同构}\leavevmode\label{def:isomorphism-of-projective-varieties}
射影簇 $X\subseteq\mathbb{P}^{m}$ 和 $Y\subseteq\mathbb{P}^{n}$ 之间的\emph{同构}, 是具有正则逆映射 $G:Y\to X$ 的正则映射 $F:X\to Y$. 若 $X$ 和 $Y$ 之间存在一个同构, 就称 $X$ 和 $Y$ \emph{同构}, 并记 $X\cong Y$.
\end{Definition}

\begin{BookFlowNote}{%
只要正则映射的复合仍是正则的 (已在 Exercise~\ref{ex:composition-of-projective-regular-maps} 中验证), 稍加思考便可看出, 同构是射影簇的集合上的等价关系.
}
\noindent 和仿射簇的情形一样, 我们把在同构下保持的结构看作射影簇的内在性质.

作为同构的第一个例子, 沿用上一节关于正则映射的讨论, 我们指出二次曲线 $\mathcal{V}(y^{2}-xz)\subseteq\mathbb{P}^{2}$ 与 $\mathbb{P}^{1}$ 同构.
\end{BookFlowNote}

\begin{Example}{$\mathcal{V}(y^{2}-xz)\cong\mathbb{P}^{1}$}\leavevmode\label{ex:projective-line-isomorphic-to-conic}
令 $X=\mathcal{V}(y^{2}-xz)\subseteq\mathbb{P}^{2}$, 考察分别在 Examples~\ref{ex:regular-map-not-polynomial-map} 和~\ref{ex:polynomial-map-from-projective-line-to-conic} 中引入的正则映射 $F:X\to\mathbb{P}^{1}$ 和 $G:\mathbb{P}^{1}\to X$:
\[
  F([a:b:c])=
  \begin{cases}
    [a:b] & [a:b:c]\notin\mathcal{V}(x,y),\\
    [b:c] & [a:b:c]\notin\mathcal{V}(y,z)
  \end{cases}
\]
以及
\[
  G([a:b])=[a^{2}:ab:b^{2}].
\]
现在证明这两个正则映射互逆. 给定点 $[a:b]\in\mathbb{P}^{1}$, 计算得
\[
  F(G([a:b]))=F([a^{2}:ab:b^{2}])
  =\left\{\begin{array}{ll}
    [a^{2}:ab] & \text{若 }a\ne0\\
    {[ab:b^{2}]} & \text{若 }b\ne0
  \end{array}\right\}=[a:b].
\]
反过来, 给定点 $[a:b:c]\in X$, 计算得
\[
  \begin{aligned}
  G(F([a:b:c]))
    &=G\left(\left\{\begin{array}{ll}
        [a:b] & \text{若 }(a,b)\ne(0,0)\\
        {[b:c]} & \text{若 }(b,c)\ne(0,0)
      \end{array}\right\}\right)\\
    &=\left\{\begin{array}{ll}
        [a^{2}:ab:b^{2}] & \text{若 }(a,b)\ne(0,0)\\
        {[b^{2}:bc:c^{2}]} & \text{若 }(b,c)\ne(0,0)
      \end{array}\right\}\\
    &=[a:b:c],
  \end{aligned}
\]
其中最后一个等式用到了对每个 $[a:b:c]\in X$ 都有 $b^{2}=ac$ 这一事实.

由于 $F$ 和 $G$ 是互逆的正则映射, 得到 $X\cong\mathbb{P}^{1}$.
\end{Example}

射影空间上的线性变换给出一类很特殊的同构. 为了建立记号, 设 $A$ 是元素在 $K$ 中的可逆 $(n+1)\times(n+1)$ 矩阵. 对每个点 $p\in\mathbb{P}^{n}$, 定义新点 $Ap\in\mathbb{P}^{n}$: 选取 $p$ 的齐次坐标, 将它们视为列向量, 再用矩阵 $A$ 左乘这个列向量. 例如, 如果
\[
  A=\begin{pmatrix}
     1&1&0\\
     1&0&2\\
     0&1&-1
  \end{pmatrix}
\]
且 $p=[a_{0}:a_{1}:a_{2}]$, 则 $Ap=[a_{0}+a_{1}:a_{0}+2a_{2}:a_{1}-a_{2}]$.

更一般地, 若记 $A=(A_{ij})$, 则上述过程给出线性映射 $F_{A}:\mathbb{P}^{n}\to\mathbb{P}^{n}$, 它由以 $A$ 各行为系数的线性多项式定义:
\[
  F_{A}=\left[\sum_{j=0}^{n}A_{0j}x_{j}:\cdots:\sum_{j=0}^{n}A_{nj}x_{j}\right].
\]
$A$ 的可逆性意味着 $F_{A}$ 在 $\mathbb{P}^{n}$ 的每一点都有定义, 而且 $F_{A}$ 的逆是正则映射
\[
  F_{A^{-1}}:\mathbb{P}^{n}\to\mathbb{P}^{n}
\]
(见 Exercise~\ref{ex:invertible-matrix-induces-projective-isomorphism}). 因此, $F_{A}$ 是 $\mathbb{P}^{n}$ 到自身的线性同构, 我们给这样由可逆矩阵产生的同构一个专门名称.

\begin{Definition}{射影等价}\leavevmode\label{def:projective-equivalence}
$\mathbb{P}^{n}$ 的一个\emph{射影等价}是形如
\[
  \begin{aligned}
    F_{A}:\mathbb{P}^{n}&\to\mathbb{P}^{n}\\
    p&\mapsto Ap,
  \end{aligned}
\]
的同构, 其中 $A$ 为元素在 $K$ 中的可逆 $(n+1)\times(n+1)$ 矩阵. 若存在射影等价 $F_{A}:\mathbb{P}^{n}\to\mathbb{P}^{n}$ 使得 $F_{A}(X)=Y$, 就称两个射影簇 $X,Y\subseteq\mathbb{P}^{n}$ \emph{射影等价}.
\end{Definition}

\BookPairedNote{
直观地看, 射影等价就是 $\mathbb{P}^{n}$ 中的线性坐标变换, 类似于向量空间中的换基. 若 $X,Y\subseteq\mathbb{P}^{n}$ 射影等价, 则它们同构. 的确, 若 $F_{A}:\mathbb{P}^{n}\to\mathbb{P}^{n}$ 是满足 $F_{A}(X)=Y$ 的射影等价, 则限制 $F_{A}:X\to Y$ 是具有正则逆映射 $F_{A^{-1}}:Y\to X$ 的正则映射.
}{
学过复分析的读者可能遇到过 $\mathbb{P}^{1}$ 的射影等价, 它们在复分析中称为\textquotedbl{}Möbius 变换\textquotedbl{}.
}
重要的是, $\mathbb{P}^{n}$ 的射影等价把射影簇映到射影簇. 具体地说, 若 $X\subseteq\mathbb{P}^{n}$ 是射影簇, 且 $F_{A}:\mathbb{P}^{n}\to\mathbb{P}^{n}$ 是射影等价, 则 $F_{A}(X)$ 也是 $\mathbb{P}^{n}$ 中的射影簇 (Exercise~\ref{ex:projective-equivalence-preserves-varieties}). 对任何正则映射 $F:\mathbb{P}^{m}\to\mathbb{P}^{n}$, 同一断言更一般地成立 (Theorem~\ref{thm:regular-images-of-projective-varieties-are-closed}), 但证明后一事实要困难得多.

下面来看 $\mathbb{P}^{2}$ 中射影等价的一个具体例子.

\begin{Example}{射影等价的二次曲线}\leavevmode\label{ex:projectively-equivalent-complex-conics}
考察两个复射影簇
\[
  X=\mathcal{V}(x^{2}+y^{2}+z^{2})\subseteq\mathbb{P}^{2}_{\mathbb{C}}
  \qquad\text{和}\qquad
  Y=\mathcal{V}(y^{2}-xz)\subseteq\mathbb{P}^{2}_{\mathbb{C}}.
\]
我们断言 $X$ 和 $Y$ 射影等价, 并将通过对定义方程作坐标变换来展示这一点. 更具体地说, 注意
\[
  x^{2}+y^{2}+z^{2}=y^{2}-(xi+z)(xi-z).
\]
从而 $[a:b:c]\in X$ 当且仅当 $[ai+c:b:ai-c]\in Y$. 换句话说, 令 $A$ 为可逆矩阵
\[
  A=\begin{pmatrix}
    i&0&1\\
    0&1&0\\
    i&0&-1
  \end{pmatrix},
\]
可知 $F_{A}(X)=Y$, 故 $X$ 和 $Y$ 射影等价.
\end{Example}

事实上, 前一例还可以用一种颇为意外的方式推广.

\begin{Proposition}{所有不可约二次曲线都射影等价}\leavevmode\label{prop:irreducible-conics-projectively-equivalent}
设 $f,g\in K[x,y,z]$ 为不可约的二次齐次多项式. 则 $\mathcal{V}(f)$ 和 $\mathcal{V}(g)$ 在 $\mathbb{P}^{2}$ 中射影等价.
\end{Proposition}
\begin{Proof}
要证明所有不可约二次曲线都射影等价, 只需证明它们都与一条特定的二次曲线射影等价. 更具体地说, 设 $f\in K[x,y,z]$ 为不可约的二次齐次多项式; 我们证明 $\mathcal{V}(f)$ 与 $\mathcal{V}(xy+xz+yz)$ 射影等价.

首先, 选取两个不同的点 $[a_{1}:b_{1}:c_{1}],[a_{2}:b_{2}:c_{2}]\in\mathcal{V}(f)$, 并考察矩阵
\[
  \begin{pmatrix}
    a_{1}&a_{2}\\
    b_{1}&b_{2}\\
    c_{1}&c_{2}
  \end{pmatrix}.
\]
由于两列互不为标量倍数, 这个矩阵的秩为二, 因而各行之间有一个在相差标量倍数意义下唯一的线性关系. 如果 $\mathcal{V}(f)$ 中任意选取的第三个点也满足这一线性关系, 则 $\mathcal{V}(f)$ 的所有点都在 $\mathbb{P}^{2}$ 中的一条直线上, 与 $f$ 的不可约性矛盾. 因此, 必有第三个点 $[a_{3}:b_{3}:c_{3}]\in\mathcal{V}(f)$ 不满足这个线性关系, 从而矩阵
\[
  A=\begin{pmatrix}
    a_{1}&a_{2}&a_{3}\\
    b_{1}&b_{2}&b_{3}\\
    c_{1}&c_{2}&c_{3}
  \end{pmatrix}
\]
的各行线性无关, 因而它可逆.

由矩阵 $A$ 定义的射影等价 $F_{A}$ 把三点 $[1:0:0],[0:1:0],[0:0:1]\in\mathbb{P}^{2}$ 映到我们在 $\mathcal{V}(f)$ 上选取的三个点. 因此, 逆射影等价 $F_{A^{-1}}$ 把 $\mathcal{V}(f)$ 等同于一条不可约二次曲线 $\mathcal{V}(g)\subseteq\mathbb{P}^{2}$, 使得 $[1:0:0],[0:1:0],[0:0:1]\in\mathcal{V}(g)$. 而二次齐次多项式在这三点处消没, 当且仅当它具有形式
\[
  g=\alpha xy+\beta xz+\gamma yz,
\]
其中 $\alpha,\beta,\gamma\in K$. 由 $g$ 的不可约性可知, $\alpha,\beta,\gamma$ 均非零. 于是, 射影等价 $[a:b:c]\mapsto[\gamma^{-1}a:\beta^{-1}b:\alpha^{-1}c]$ 把 $\mathcal{V}(g)$ 等同于 $\mathcal{V}(xy+xz+yz)$. 将以上射影等价复合, 可知 $\mathcal{V}(f)$ 与 $\mathcal{V}(xy+xz+yz)$ 射影等价.
\end{Proof}

\begin{Corollary}{所有不可约二次曲线都同构于 $\mathbb{P}^{1}$}\leavevmode\label{cor:irreducible-conics-isomorphic-to-projective-line}
设 $f\in K[x,y,z]$ 为不可约的二次齐次多项式. 则
\[
  \mathcal{V}(f)\cong\mathbb{P}^{1}.
\]
\end{Corollary}
\begin{Proof}
由 Proposition~\ref{prop:irreducible-conics-projectively-equivalent}, 所有不可约二次曲线都射影等价, 因而同构. 再由 Example~\ref{ex:projective-line-isomorphic-to-conic}, 不可约二次曲线 $\mathcal{V}(y^{2}-xz)$ 与 $\mathbb{P}^{1}$ 同构. 因此 $\mathbb{P}^{2}$ 中所有不可约二次曲线都同构于 $\mathbb{P}^{1}$.
\end{Proof}

Corollary~\ref{cor:irreducible-conics-isomorphic-to-projective-line} 表明, $\mathbb{P}^{2}$ 中不可约二次曲线的内在性质相当简单, 因为它们都同构于 $\mathbb{P}^{1}$. 但要提醒读者, 不要因此以为 $\mathbb{P}^{2}$ 中更高次曲线的内在性质也同样简单. 事实上, 只需把次数提高一次, 就可以证明 $\mathbb{P}^{2}$ 中的不可约三次曲线有无穷多个不同的同构类, 它们都不与 $\mathbb{P}^{1}$ 同构 (在 Example~\ref{ex:irreducible-projective-cubic-not-isomorphic-to-line} 中, 我们给出一个不与 $\mathbb{P}^{1}$ 同构的不可约三次曲线的具体例子). $\mathbb{P}^{2}$ 中的三次曲线是研究椭圆曲线的基础. 椭圆曲线是一个引人入胜的数学分支, 许多研究者将整个职业生涯投入其中, 不过本书不会涉及这一主题.

关于同构的讨论自然引出一个问题: 如何研究射影簇的内在性质? 在仿射情形下, 我们引入了仿射簇的坐标环作为这方面的关键工具: 两个仿射簇同构当且仅当它们的坐标环同构, 从而可以通过研究坐标环的代数结构来研究仿射簇的内在性质. 受仿射情形的启发, 不妨问是否存在某个代数对象, 能刻画射影簇的内在性质. 沿用仿射情形中的发展, 很自然地可以引入射影簇 $X\subseteq\mathbb{P}^{n}$ 的\emph{齐次坐标环}, 定义为商环
\[
  \frac{K[x_{0},\ldots,x_{n}]}{\mathcal{I}(X)}.
\]
遗憾的是, 与仿射情形不同, 齐次坐标环在同构下并不保持, 如下一例所示.

\begin{Example}{同构与齐次坐标环}\leavevmode\label{ex:isomorphism-does-not-preserve-homogeneous-coordinate-ring}
Example~\ref{ex:projective-line-isomorphic-to-conic} 表明 $X=\mathcal{V}(y^{2}-xz)\subseteq\mathbb{P}^{2}$ 与 $\mathbb{P}^{1}$ 同构. 不过, 这两个射影簇的齐次坐标环分别为
\[
  \frac{K[x,y,z]}{\langle y^{2}-xz\rangle}
  \qquad\text{和}\qquad K[s,t],
\]
而这两个环并不同构. 特别是, 后者是唯一分解整环, 但方程 $y^{2}=xz$ 蕴含前者不是.

也可以用另一种方式看出这两个环并不同构: 如果它们同构, 则由仿射簇情形下代数与几何的等价性, 仿射簇 $\mathcal{V}_{\mathbb{A}}(y^{2}-xz)$ 和 $\mathbb{A}^{2}$ 应当同构. 但它们并不同构: 前者在原点有一个奇点, 而后者是非奇异的.
\end{Example}

\BookPairedNoteLeft{
齐次坐标环在射影代数几何中起重要作用, 但不是本书论述的重点.
}{
尽管前一例说明齐次坐标环不是射影簇的完全代数不变量, 仍有一种可能: 每个射影簇或许还对应另一个环, 能够编码它的全部内在结构.
}
事实并非如此. 正如下一章将看到的, 射影簇的内在性质不是由单独一个环决定的, 而是由一族环决定的. 每个环记录射影簇的局部几何信息, 整体几何则由这些局部环之间的同态决定.

不过, 暂且不要走得太远. 在本章余下部分, 我们将介绍 Veronese, Segre 和 Plücker 这三族经典映射, 继续加深对射影簇之间映射的实际理解.

\BookExerciseHeading
\begin{BookExercises}
  \item\label{ex:twisted-cubic-isomorphic-to-projective-line} 令 $X=\mathcal{V}(x_{0}x_{2}-x_{1}^{2},x_{0}x_{3}-x_{1}x_{2},x_{1}x_{3}-x_{2}^{2})\subseteq\mathbb{P}^{3}$ 为扭三次曲线. 考察映射
  \[
    \begin{aligned}
      F:\mathbb{P}^{1}&\to X\\
      [a:b]&\mapsto[a^{3}:a^{2}b:ab^{2}:b^{3}].
    \end{aligned}
  \]
  通过构造正则逆映射, 证明 $F$ 是同构.
  \item\label{ex:linear-variety-of-independent-polynomials-is-projective-space} 设 $\ell_{1},\ldots,\ell_{k}\in K[x_{0},\ldots,x_{n}]$ 为线性无关的齐次线性多项式. 证明 $\mathcal{V}(\ell_{1},\ldots,\ell_{k})\cong\mathbb{P}^{n-k}$.
  \item\label{ex:invertible-matrix-induces-projective-isomorphism} 设 $A$ 是元素在 $K$ 中的可逆 $(n+1)\times(n+1)$ 矩阵.
  \begin{enumerate}
    \item 证明 $F_{A}:\mathbb{P}^{n}\to\mathbb{P}^{n}$ 是正则映射.
    \item 证明 $F_{A^{-1}}$ 是 $F_{A}$ 的逆.
  \end{enumerate}
  \item\label{ex:projective-equivalence-preserves-varieties} 设 $A$ 是元素在 $K$ 中的可逆 $(n+1)\times(n+1)$ 矩阵.
  \begin{enumerate}
    \item 令 $X=\mathcal{V}(f_{1},\ldots,f_{k})\in\mathbb{P}^{n}$. 证明
    \[
      F_{A}(X)=\mathcal{V}(f_{1}\circ A^{-1},\ldots,f_{k}\circ A^{-1}).
    \]
    (这里, 自然地把每个 $f_{i}$ 看作函数 $f_{i}:K^{n+1}\to K$, 把 $A^{-1}$ 看作函数 $A^{-1}:K^{n+1}\to K^{n+1}$.) 特别地, 若 $X$ 为射影簇, 则 $F_{A}(X)$ 也是射影簇.
    \item 设 $X,Y\subseteq\mathbb{P}^{n}$ 射影等价. 证明 $X$ 不可约当且仅当 $Y$ 不可约.
  \end{enumerate}
  \item\label{ex:projective-equivalence-preserves-homogeneous-coordinate-ring} 设 $X,Y\subseteq\mathbb{P}^{n}$ 为射影等价的两个射影簇. 证明
  \[
    \frac{K[x_{0},\ldots,x_{n}]}{\mathcal{I}(X)}
    \cong
    \frac{K[x_{0},\ldots,x_{n}]}{\mathcal{I}(Y)}.
  \]
  因此, 虽然齐次坐标环在同构下不保持, 但在射影等价下保持.
  \item\label{ex:normal-form-and-finite-fibres-of-projective-linear-map} 设 $X\subseteq\mathbb{P}^{n}$ 为射影簇, 设 $\ell_{0},\ldots,\ell_{k}\in K[x_{0},\ldots,x_{n}]$ 为齐次线性多项式, 且 $\mathcal{V}(\ell_{0},\ldots,\ell_{k})\cap X=\varnothing$. 考察线性映射
  \[
    \begin{aligned}
      L:X&\to\mathbb{P}^{k}\\
      a&\mapsto[\ell_{0}(a):\cdots:\ell_{k}(a)].
    \end{aligned}
  \]
  \begin{enumerate}
    \item 证明存在一对射影等价 $F_{A}:\mathbb{P}^{n}\to\mathbb{P}^{n}$ 和 $F_{B}:\mathbb{P}^{k}\to\mathbb{P}^{k}$, 使对所有 $a=[a_{0}:\cdots:a_{n}]\in X$, 都有
    \[
      F_{B}(L(F_{A}(a)))=[a_{0}:\cdots:a_{j}],
    \]
    其中 $j\le k$. (提示: 对以 $\ell_{1},\ldots,\ell_{k}$ 的系数为各行的 $k\times n$ 矩阵作行, 列运算, 可以将其化为形如
    \[
      \begin{pmatrix}
        I&0\\
        0&0
      \end{pmatrix}
    \]
    的矩阵, 其中 $I$ 表示某个 $j\le k$ 对应的 $j\times j$ 单位矩阵, 每个 $0$ 表示相应大小的零矩阵.)
    \item 利用 (a) 证明 $L$ 是有限对一的. (提示: 见 Exercise~\ref{ex:projective-linear-projection-is-finite-to-one}.)
  \end{enumerate}
\end{BookExercises}
